\documentclass[a4paper,11pt]{article}

% ── Packages ──
\usepackage[utf8]{inputenc}
\usepackage[T1]{fontenc}
\usepackage[italian]{babel}
\usepackage{amsmath,amssymb,amsthm,mathtools}
\usepackage{bm}
\usepackage{geometry}
\geometry{margin=2.5cm}
\usepackage{hyperref}
\hypersetup{colorlinks=true, linkcolor=blue!70!black, citecolor=green!50!black, urlcolor=blue!60!black}
\usepackage{enumitem}
\usepackage{booktabs}
\usepackage{tabularx}
\usepackage{xcolor}
\usepackage{tcolorbox}
\usepackage{tikz}
\usepackage{float}

% ── Colored boxes ──
\newtcolorbox{phasebox}[1]{colback=blue!5!white, colframe=blue!60!black, title={\textbf{#1}}, fonttitle=\large}
\newtcolorbox{tipbox}{colback=green!5!white, colframe=green!50!black, title={\textbf{Suggerimento}}}
\newtcolorbox{warnbox}{colback=red!5!white, colframe=red!50!black, title={\textbf{Attenzione}}}
\newtcolorbox{splitbox}[1]{colback=orange!5!white, colframe=orange!60!black, title={\textbf{Divisione del lavoro: #1}}}

% ── Notation ──
\newcommand{\R}{\mathbb{R}}
\newcommand{\norm}[1]{\left\lVert #1 \right\rVert}
\DeclareMathOperator{\diag}{diag}
\DeclareMathOperator{\sign}{sign}
\newcommand{\argmin}{\operatornamewithlimits{arg\,min}}

% ══════════════════════════════════════════════════════════════
\title{\textbf{Guida allo Sviluppo --- Progetto 25}\\[6pt]
\large Piano di lavoro, struttura del codice e suddivisione tra due persone\\[8pt]
\normalsize Computational Mathematics for Learning and Data Analysis (646AA)}
\author{}
\date{A.A.\ 2025/26}

\begin{document}
\maketitle
\tableofcontents
\newpage

% ══════════════════════════════════════════════════════════════
\section{Panoramica del progetto}
\label{sec:overview}

Il Progetto~25 richiede di:
\begin{enumerate}
  \item \textbf{(M)} Costruire un modello \emph{Extreme Learning Machine} (ELM): una rete neurale a uno strato nascosto con pesi casuali fissi e strato di output appreso tramite un problema LASSO ($L_1$-regolarizzato).
  \item \textbf{(A1)} Implementare l'algoritmo \emph{Iteratively Reweighted Least Squares} (IRLS) per risolvere il LASSO.
  \item \textbf{(A2)} Implementare un algoritmo della classe \emph{Deflected Subgradient Methods} per risolvere lo stesso problema.
  \item Confrontare i due algoritmi in termini di convergenza, accuratezza e efficienza.
\end{enumerate}

\begin{warnbox}
\textbf{Vincoli importanti (dal testo del progetto):}
\begin{itemize}[nosep]
  \item ``No off-the-shelf solvers allowed'' per la risoluzione del problema LASSO. Si possono usare librerie per operazioni di base (es.\ \texttt{numpy.linalg.solve} per i sotto-sistemi lineari dell'IRLS).
  \item Il report deve seguire la struttura indicata nelle istruzioni del corso.
  \item Il codice deve essere ben organizzato, non tutto in un singolo notebook.
\end{itemize}
\end{warnbox}

% ──────────────────────────────────────────────────────────────
\subsection{Linguaggio e strumenti consigliati}
\label{subsec:tools}

\begin{center}
\begin{tabular}{ll}
\toprule
\textbf{Aspetto} & \textbf{Consiglio} \\
\midrule
Linguaggio & Python (con NumPy/SciPy) \\
Algebra lineare & \texttt{numpy.linalg}, \texttt{scipy.linalg} \\
Grafici & \texttt{matplotlib} \\
Dati & \texttt{scikit-learn} (solo per dataset e confronto) \\
Versioning & \texttt{git} (repo privato!) \\
Report & \LaTeX \\
\bottomrule
\end{tabular}
\end{center}

% ══════════════════════════════════════════════════════════════
\section{Struttura dei file}
\label{sec:structure}

\begin{verbatim}
cm_project/
+-- src/
|   +-- elm.py                  # Classe ELM (modello M)
|   +-- irls.py                 # Algoritmo A1 (IRLS)
|   +-- deflected_subgradient.py # Algoritmo A2
|   +-- lasso_utils.py          # Funzioni comuni (f, grad, ...)
|   +-- linear_solvers.py       # Cholesky, CG per IRLS
|   +-- data_generation.py      # Generazione/caricamento dati
+-- experiments/
|   +-- experiment_convergence.py  # Test convergenza e plot
|   +-- experiment_comparison.py   # Confronto A1 vs A2
|   +-- experiment_scalability.py  # Scalabilita'
|   +-- experiment_params.py       # Studio parametri
+-- notebooks/
|   +-- workflow.ipynb           # Notebook riassuntivo
+-- results/
|   +-- figures/                 # Grafici generati
|   +-- tables/                  # Tabelle CSV con risultati
+-- theory/
|   +-- guida_teorica.tex
|   +-- guida_sviluppo.tex
+-- report/
|   +-- report.tex               # Report finale
+-- README.md
+-- requirements.txt
\end{verbatim}

% ══════════════════════════════════════════════════════════════
\section{Piano di lavoro suddiviso in fasi}
\label{sec:phases}

Il progetto si divide naturalmente in \textbf{6 fasi}. Per ciascuna fase indichiamo i task, le dipendenze, e una proposta di suddivisione del lavoro tra \textbf{Persona A} e \textbf{Persona B}.

% ──────────────────────────────────────────────────────────────
\begin{phasebox}{Fase 0: Setup e comprensione (1 settimana)}

\textbf{Obiettivi:}
\begin{itemize}[nosep]
  \item Comprendere a fondo il testo del progetto e la teoria sottostante.
  \item Impostare l'ambiente di sviluppo e la struttura dei file.
  \item Scrivere la sezione teorica del report (``Setting the stage'').
\end{itemize}

\textbf{Task:}
\begin{enumerate}[nosep]
  \item Leggere il testo del progetto e i riferimenti citati (soprattutto d'Antonio--Frangioni 2009 per i deflected subgradient).
  \item Creare il repository git e la struttura dei file.
  \item Implementare le funzioni di base in \texttt{lasso\_utils.py}:
  \begin{itemize}[nosep]
    \item $f(w) = \norm{Xw - y}_2^2 + \lambda\norm{w}_1$
    \item Gradiente della parte smooth: $\nabla g(w) = 2X^T(Xw - y)$
    \item Subgradiente di $f$: $\partial f(w)$
    \item Metrica di gap/residuo
  \end{itemize}
  \item Scrivere le prime sezioni del report (\S 4.2 delle istruzioni).
\end{enumerate}

\begin{splitbox}{Fase 0}
\begin{tabularx}{\textwidth}{lX}
\textbf{Persona A:} & Setup repo, struttura file, implementazione \texttt{lasso\_utils.py}, test unitari. \\
\textbf{Persona B:} & Studio teoria, prima bozza sezione teorica del report (definizioni, formulazione del problema, condizioni di ottimalità). \\
\textbf{Insieme:} & Discussione e comprensione del progetto.
\end{tabularx}
\end{splitbox}
\end{phasebox}

% ──────────────────────────────────────────────────────────────
\begin{phasebox}{Fase 1: Modello ELM e generazione dati (1 settimana)}

\textbf{Obiettivi:}
\begin{itemize}[nosep]
  \item Implementare la classe ELM completa.
  \item Preparare dataset sintetici e reali per i test.
\end{itemize}

\textbf{Task:}
\begin{enumerate}[nosep]
  \item \textbf{Classe ELM} (\texttt{elm.py}):
  \begin{itemize}[nosep]
    \item \texttt{\_\_init\_\_(d, p, activation='sigmoid')}: genera $W_1$ casuale.
    \item \texttt{transform(X\_raw)}: calcola $\sigma(W_1 X_{\text{raw}}^T)$, restituisce la matrice delle feature nascoste.
    \item \texttt{fit(X\_hidden, y, solver)}: chiama il solver scelto (IRLS o deflected subgradient).
    \item \texttt{predict(X\_raw)}: applica il modello completo.
  \end{itemize}
  \item \textbf{Generazione dati} (\texttt{data\_generation.py}):
  \begin{itemize}[nosep]
    \item Dataset sintetici: generare $w^*$ sparso, $X$ casuale, $y = Xw^* + \text{noise}$.
    \item Variare: dimensione ($n$, $m$), sparsità di $w^*$, livello di rumore.
    \item Opzionale: dataset reali (es.\ da UCI, sklearn).
  \end{itemize}
\end{enumerate}

\begin{tipbox}
Per la generazione dei dati sintetici, è utile conoscere la soluzione ottima $w^*$ (o almeno $f^*$) per poter valutare il gap. Un modo: generare $w^*$ sparso, calcolare $y = Xw^*$, e usare un solutore esterno (es.\ \texttt{sklearn.linear\_model.Lasso}) per ottenere un riferimento $f^*_{\text{ref}}$.
\end{tipbox}

\begin{splitbox}{Fase 1}
\begin{tabularx}{\textwidth}{lX}
\textbf{Persona A:} & Classe ELM, scelta della funzione di attivazione, test. \\
\textbf{Persona B:} & Generazione dati sintetici, eventuale importazione dati reali, utilità per il calcolo di $f^*_{\text{ref}}$. \\
\end{tabularx}
\end{splitbox}
\end{phasebox}

% ──────────────────────────────────────────────────────────────
\begin{phasebox}{Fase 2: Implementazione di A1 --- IRLS (2 settimane)}

\textbf{Obiettivi:}
\begin{itemize}[nosep]
  \item Implementare l'algoritmo IRLS completo.
  \item Implementare i solutori lineari necessari (Cholesky, CG).
  \item Testare su piccole istanze.
\end{itemize}

\textbf{Task dettagliati:}

\paragraph{2a. Solutori lineari (\texttt{linear\_solvers.py}):}
\begin{enumerate}[nosep]
  \item \textbf{Risoluzione via equazioni normali + Cholesky:}
  \begin{itemize}[nosep]
    \item Input: $Q \succ 0$, $b$. Output: $w$ tale che $Qw = b$.
    \item Si può usare \texttt{scipy.linalg.cho\_factor} / \texttt{cho\_solve}, oppure implementare la propria fattorizzazione.
    \item \emph{Nota:} \texttt{numpy.linalg.solve} è permesso come sotto-passo.
  \end{itemize}
  \item \textbf{Gradiente Coniugato (opzionale ma consigliato):}
  \begin{itemize}[nosep]
    \item Implementare CG come descritto nelle slide di NLA, Lecture~5.
    \item Vantaggioso per grandi $n$ poiché evita di formare $X^TX$ esplicitamente.
    \item Operazione dominante: prodotto matrice-vettore $Q v$.
  \end{itemize}
\end{enumerate}

\paragraph{2b. IRLS (\texttt{irls.py}):}
\begin{enumerate}[nosep]
  \item Implementare l'Algoritmo~1 della guida teorica.
  \item Parametri: $\varepsilon_{\text{thr}}$, $\varepsilon_{\text{stop}}$, $k_{\max}$, tipo di solver lineare.
  \item Registrare ad ogni iterazione: $f(w_k)$, $\norm{w_{k+1}-w_k}$, tempo CPU.
  \item \textbf{Inizializzazione:} $w_0 = (X^TX)^{-1}X^Ty$ (minimi quadrati standard) è una buona scelta.
\end{enumerate}

\paragraph{2c. Test e debug:}
\begin{enumerate}[nosep]
  \item Verificare su un piccolo esempio ($n=5$, $m=20$) che IRLS converge.
  \item Confrontare la soluzione con \texttt{sklearn.linear\_model.Lasso} per validare.
  \item Verificare che $0 \in \partial f(w^*)$ alla convergenza (controllare le condizioni di ottimalità componentwise).
\end{enumerate}

\begin{splitbox}{Fase 2}
\begin{tabularx}{\textwidth}{lX}
\textbf{Persona A:} & Solutori lineari (Cholesky + CG), test unitari dei solutori. \\
\textbf{Persona B:} & IRLS completo, integrazione con solutori, test di convergenza e confronto con solutori di riferimento. \\
\textbf{Insieme:} & Debug, revisione codice reciproca.
\end{tabularx}
\end{splitbox}

\begin{warnbox}
Aspetti critici dell'IRLS:
\begin{itemize}[nosep]
  \item Il threshold $\varepsilon_{\text{thr}}$ influenza pesantemente la convergenza e la qualità della soluzione. Testare diversi valori ($10^{-6}$, $10^{-8}$, $10^{-10}$).
  \item Se $X$ non ha rango pieno, $X^TX$ è singolare: il termine $2\lambda W_k^TW_k$ regolarizza, ma attenzione al condizionamento.
  \item Pre-calcolare $X^TX$ e $X^Ty$ una volta per tutte (\emph{non} ricalcolarli ad ogni iterazione).
\end{itemize}
\end{warnbox}
\end{phasebox}

% ──────────────────────────────────────────────────────────────
\begin{phasebox}{Fase 3: Implementazione di A2 --- Deflected Subgradient (2 settimane)}

\textbf{Obiettivi:}
\begin{itemize}[nosep]
  \item Implementare l'algoritmo deflected subgradient con target level.
  \item Calibrare i parametri.
  \item Testare su piccole istanze e confrontare con A1.
\end{itemize}

\textbf{Task dettagliati:}

\paragraph{3a. Subgradiente puro (baseline):}
\begin{enumerate}[nosep]
  \item Implementare il subgradiente standard con stepsize decrescente ($\alpha^i = c/i$) come baseline di confronto.
  \item Implementare il Polyak stepsize (usando $f^*_{\text{ref}}$ dal solutore di riferimento).
  \item Implementare il target level stepsize.
\end{enumerate}

\paragraph{3b. Deflected subgradient (\texttt{deflected\_subgradient.py}):}
\begin{enumerate}[nosep]
  \item Implementare l'Algoritmo~2 della guida teorica.
  \item Implementare le regole di deflection:
  \begin{itemize}[nosep]
    \item Deflection ottimale \eqref{eq:gamma_opt}: formula chiusa come proiezione del minimo quadratico su $[0,1]$.
    \item Stepsize-restricted con $\beta^i \leq \gamma^i$.
    \item Deflection-restricted.
  \end{itemize}
  \item Registrare ad ogni iterazione: $f(w_i)$, $\bar{f}^i$, target level, $\norm{d^i}$, tempo CPU.
\end{enumerate}

\paragraph{3c. Formula chiusa per $\gamma^i$:}
Il problema $\min_{\gamma \in [0,1]} \norm{\gamma g + (1-\gamma) d^{-}}^2$ è un problema quadratico univariato in $\gamma$:
\[
  \phi(\gamma) = \gamma^2\norm{g}^2 + 2\gamma(1-\gamma)\langle g, d^{-}\rangle + (1-\gamma)^2\norm{d^{-}}^2.
\]
Il minimo non vincolato è:
\[
  \gamma^* = \frac{\norm{d^{-}}^2 - \langle g, d^{-}\rangle}{\norm{g}^2 - 2\langle g, d^{-}\rangle + \norm{d^{-}}^2} = \frac{\norm{d^{-}}^2 - \langle g, d^{-}\rangle}{\norm{g - d^{-}}^2}.
\]
Poi si proietta su $[0, 1]$: $\gamma^i = \min(1, \max(0, \gamma^*))$.

\paragraph{3d. Test e calibrazione:}
\begin{enumerate}[nosep]
  \item Testare su piccole istanze con $f^*$ noto.
  \item Studiare l'effetto dei parametri $\delta_0$, $\rho$, $R$, $\beta$ sulla convergenza.
  \item Confrontare subgradiente puro vs.\ deflected subgradient.
\end{enumerate}

\begin{splitbox}{Fase 3}
\begin{tabularx}{\textwidth}{lX}
\textbf{Persona A:} & Subgradiente puro (baseline), implementazione delle diverse stepsize rules. \\
\textbf{Persona B:} & Deflected subgradient, regole di deflection, formula chiusa per $\gamma$. \\
\textbf{Insieme:} & Calibrazione parametri, confronti iniziali A1 vs A2.
\end{tabularx}
\end{splitbox}

\begin{tipbox}
Suggerimenti per il debugging del deflected subgradient:
\begin{itemize}[nosep]
  \item Verificare che il record value $\bar{f}^i$ sia sempre $\leq f(w_0)$.
  \item Verificare che la direzione $d^i$ sia una combinazione convessa di $g^i$ e $d^{i-1}$ (se $\gamma^i \in [0,1]$).
  \item Plottare l'evoluzione di $\gamma^i$ e $\alpha^i$ nel tempo: se $\gamma^i \to 1$ rapidamente, la deflection non sta aiutando.
  \item Confrontare con il Polyak stepsize (che richiede $f^*$) per verificare che il target level sta funzionando.
\end{itemize}
\end{tipbox}
\end{phasebox}

% ──────────────────────────────────────────────────────────────
\begin{phasebox}{Fase 4: Esperimenti e analisi (2 settimane)}

\textbf{Obiettivi:}
\begin{itemize}[nosep]
  \item Condurre esperimenti sistematici.
  \item Produrre grafici e tabelle significative.
  \item Analizzare e commentare i risultati.
\end{itemize}

\textbf{Esperimenti da condurre:}

\paragraph{4a. Convergenza:}
\begin{itemize}[nosep]
  \item Plottare $f(w_k) - f^*$ (o $\bar{f}^k - f^*$) vs.\ iterazioni in scala log-lineare.
  \item Plottare $f(w_k) - f^*$ vs.\ tempo CPU.
  \item Verificare il tasso di convergenza (lineare per IRLS, sublineare per deflected subgradient).
  \item \textbf{Scala logaritmica:} essenziale per osservare la convergenza lineare (appare come retta nel grafico log-lineare).
\end{itemize}

\paragraph{4b. Scalabilità:}
\begin{itemize}[nosep]
  \item Variare $n$ (dimensione del problema): $n = 50, 100, 500, 1000, 5000$.
  \item Variare $m$ (numero di campioni): $m = 2n, 5n, 10n$.
  \item Misurare tempo CPU totale e per iterazione.
  \item Identificare il massimo $n$ risolvibile in tempo ragionevole.
\end{itemize}

\paragraph{4c. Effetto dei parametri:}
\begin{itemize}[nosep]
  \item Per IRLS: effetto di $\varepsilon_{\text{thr}}$ ($10^{-4}$, $10^{-6}$, $10^{-8}$, $10^{-10}$, $10^{-12}$).
  \item Per A2: effetto di $\delta_0$, $\rho$, $R$, tipo di deflection.
  \item Effetto di $\lambda$ (regolarizzazione) sulla sparsità e sulla convergenza.
  \item Per ELM: effetto di $p$ (numero di neuroni nascosti) e della funzione di attivazione.
\end{itemize}

\paragraph{4d. Confronto A1 vs A2:}
\begin{itemize}[nosep]
  \item Confronto iterazioni necessarie per raggiungere $f(w) - f^* \leq \varepsilon$ per diversi $\varepsilon$.
  \item Confronto tempo CPU per raggiungere la stessa accuratezza.
  \item Confronto su diverse dimensioni del problema.
  \item A1 ha iterazioni costose ma poche; A2 ha iterazioni economiche ma molte: chi ``vince''?
\end{itemize}

\paragraph{4e. Validazione:}
\begin{itemize}[nosep]
  \item Confronto con \texttt{sklearn.linear\_model.Lasso} per verificare la correttezza del valore $f^*$.
  \item Confronto della soluzione $w^*$ (norma della differenza).
  \item Verifica delle condizioni di ottimalità \eqref{eq:opt_cond}.
\end{itemize}

\begin{splitbox}{Fase 4}
\begin{tabularx}{\textwidth}{lX}
\textbf{Persona A:} & Esperimenti di convergenza e scalabilità per A1 (IRLS). Confronto solutori lineari (Cholesky vs CG). \\
\textbf{Persona B:} & Esperimenti di convergenza e calibrazione parametri per A2 (deflected subgradient). \\
\textbf{Insieme:} & Confronto A1 vs A2, grafici finali, analisi dei risultati.
\end{tabularx}
\end{splitbox}

\begin{tipbox}
Per i grafici, seguire le indicazioni del corso (\S 4.6 delle istruzioni):
\begin{itemize}[nosep]
  \item Usare scale logaritmiche per residui e gap.
  \item La convergenza lineare appare come retta in un grafico log-lineare.
  \item Pochi grafici significativi sono meglio di molti grafici illeggibili.
  \item Ogni grafico deve avere un \emph{purpose} chiaro.
\end{itemize}
\end{tipbox}
\end{phasebox}

% ──────────────────────────────────────────────────────────────
\begin{phasebox}{Fase 5: Report e consegna (1--2 settimane)}

\textbf{Obiettivi:}
\begin{itemize}[nosep]
  \item Completare il report seguendo la struttura richiesta.
  \item Preparare il pacchetto di consegna (codice + dati + risultati + report).
\end{itemize}

\textbf{Struttura del report} (dalle istruzioni del corso, \S 4):
\begin{enumerate}
  \item \textbf{Frontespizio:} numero del gruppo, numero del progetto, testo esatto della richiesta.
  \item \textbf{Descrizione del problema e degli algoritmi (\S 4.2):}
  \begin{itemize}[nosep]
    \item Formulazione del problema (ELM + LASSO).
    \item Breve richiamo della teoria necessaria (non ricopiare le slides!).
    \item Descrizione degli algoritmi con i dettagli matematici dell'adattamento al problema specifico.
    \item Motivazione delle scelte fatte.
  \end{itemize}
  \item \textbf{Proprietà attese (\S 4.3):}
  \begin{itemize}[nosep]
    \item Risultati di convergenza: quali teoremi si applicano? Le ipotesi sono soddisfatte?
    \item Velocità di convergenza attesa.
    \item Complessità computazionale.
  \end{itemize}
  \item \textbf{Setup sperimentale (\S 4.5):}
  \begin{itemize}[nosep]
    \item Descrizione dei dati (sintetici e/o reali).
    \item Scelta dei parametri algoritmici.
  \end{itemize}
  \item \textbf{Risultati sperimentali (\S 4.6):}
  \begin{itemize}[nosep]
    \item Grafici e tabelle con commenti.
    \item Confronto tra gli algoritmi.
    \item Confronto con solutori di riferimento.
  \end{itemize}
\end{enumerate}

\begin{splitbox}{Fase 5}
\begin{tabularx}{\textwidth}{lX}
\textbf{Persona A:} & Sezioni del report su A1, risultati sperimentali di IRLS, revisione finale codice. \\
\textbf{Persona B:} & Sezioni del report su A2, risultati sperimentali del deflected subgradient, README e documentazione. \\
\textbf{Insieme:} & Sezioni introduttive e conclusive, revisione incrociata completa, preparazione del pacchetto di consegna.
\end{tabularx}
\end{splitbox}

\begin{warnbox}
Dalle istruzioni del corso:
\begin{itemize}[nosep]
  \item \textbf{Non} copiare testo da slides, libri, internet. Scrivere con parole proprie.
  \item \textbf{Non} scrivere lunghi paragrafi su argomenti non rilevanti (es.\ applicazioni pratiche del modello, teoria del ML). Concentrarsi sulla parte \emph{matematica e algoritmica}.
  \item Citare sempre le fonti: ``Questo teorema è da [X, Theorem Y]''.
  \item Si consiglia di inviare il report \emph{parziale} ai docenti per feedback prima di completare tutto.
\end{itemize}
\end{warnbox}
\end{phasebox}

% ══════════════════════════════════════════════════════════════
\section{Riepilogo della divisione del lavoro}
\label{sec:summary_split}

\begin{center}
\renewcommand{\arraystretch}{1.3}
\begin{tabularx}{\textwidth}{|c|X|X|c|}
\hline
\textbf{Fase} & \textbf{Persona A} & \textbf{Persona B} & \textbf{Durata} \\
\hline\hline
0 & Setup, utilities, test & Teoria, bozza report & 1 sett.\ \\
\hline
1 & Classe ELM & Dati sintetici e reali & 1 sett.\ \\
\hline
2 & Solutori lineari (Cholesky, CG) & IRLS completo & 2 sett.\ \\
\hline
3 & Subgradiente puro (baseline) & Deflected subgradient & 2 sett.\ \\
\hline
4 & Esperimenti A1 + scalabilità & Esperimenti A2 + parametri & 2 sett.\ \\
\hline
5 & Report A1 + revisione & Report A2 + README & 1--2 sett.\ \\
\hline
\multicolumn{3}{|r|}{\textbf{Totale stimato}} & \textbf{9--10 sett.} \\
\hline
\end{tabularx}
\end{center}

\begin{tipbox}
La suddivisione è indicativa. L'aspetto cruciale è che entrambe le persone comprendano \emph{entrambi} gli algoritmi, poiché l'esame orale può vertere su qualsiasi parte del progetto.
\end{tipbox}

% ══════════════════════════════════════════════════════════════
\section{Dettagli implementativi}
\label{sec:implementation}

\subsection{Funzioni comuni (\texttt{lasso\_utils.py})}
\label{subsec:lasso_utils}

\begin{verbatim}
import numpy as np

def f_lasso(X, y, w, lam):
    """Funzione obiettivo: ||Xw - y||^2 + lambda * ||w||_1"""
    residual = X @ w - y
    return np.dot(residual, residual) + lam * np.sum(np.abs(w))

def grad_smooth(X, y, w):
    """Gradiente della parte smooth: 2 X^T (Xw - y)"""
    return 2 * X.T @ (X @ w - y)

def subgradient_f(X, y, w, lam):
    """Un subgradiente di f in w"""
    g_smooth = grad_smooth(X, y, w)
    # Per w_i != 0, sign(w_i); per w_i = 0, scegli 0
    s = np.sign(w)  # np.sign(0) = 0
    return g_smooth + lam * s

def optimality_gap(X, y, w, lam, f_star):
    """Gap rispetto all'ottimo (se noto)"""
    return f_lasso(X, y, w, lam) - f_star
\end{verbatim}

\subsection{Pre-calcolo di \texorpdfstring{$X^TX$}{XTX} e \texorpdfstring{$X^Ty$}{XTy}}
Per l'IRLS, pre-calcolare:
\begin{verbatim}
XtX = X.T @ X      # O(mn^2), fatto una sola volta
Xty = X.T @ y      # O(mn), fatto una sola volta
\end{verbatim}
Ad ogni iterazione IRLS, il sistema diventa $(X^TX + 2\lambda W_k^2) w = X^Ty$, dove $W_k^2$ è diagonale e va solo aggiunta alla diagonale di $X^TX$.

\subsection{Criteri di arresto}

\paragraph{Per IRLS:}
\begin{itemize}[nosep]
  \item Convergenza relativa: $\norm{w_{k+1} - w_k} / \max(1, \norm{w_k}) < \varepsilon$.
  \item Numero massimo di iterazioni.
  \item (Opzionale) Variazione relativa di $f$: $|f(w_{k+1}) - f(w_k)| / |f(w_k)| < \varepsilon$.
\end{itemize}

\paragraph{Per Deflected Subgradient:}
\begin{itemize}[nosep]
  \item Numero massimo di iterazioni (principale).
  \item Gap stimato $\bar{f}^i - (f_{\text{ref}} - \delta) < \varepsilon$ (poco affidabile in pratica).
  \item Tempo CPU massimo.
\end{itemize}

\subsection{Suggerimenti numerici}

\begin{enumerate}
  \item \textbf{Evitare la formazione esplicita di $X^TX$ se possibile.} Per CG, il prodotto $(X^TX + D)v$ si può calcolare come $X^T(Xv) + Dv$ senza mai formare $X^TX$.

  \item \textbf{Seed casuale fisso.} Per la riproducibilità, fissare \texttt{np.random.seed(42)} (o usare \texttt{np.random.RandomState}).

  \item \textbf{Misura del tempo.} Usare \texttt{time.perf\_counter()} per misurare i tempi CPU. Non includere il tempo di generazione dei dati o di calcolo delle metriche nel tempo dell'algoritmo.

  \item \textbf{Warm starting.} Usare la soluzione dell'esecuzione precedente come punto iniziale quando si cambia un parametro (es.\ $\lambda$).

  \item \textbf{Normalizzazione dei dati.} Normalizzare le colonne di $X$ (media 0, varianza 1) può migliorare il condizionamento e la convergenza.
\end{enumerate}

% ══════════════════════════════════════════════════════════════
\section{Checklist finale}
\label{sec:checklist}

Prima della consegna, verificare:
\begin{enumerate}[label=\fbox{\phantom{X}}, leftmargin=2cm]
  \item Il codice gira senza errori partendo da zero (testare su una macchina ``pulita'').
  \item Il README spiega come eseguire il codice e riprodurre i risultati.
  \item I grafici nel report sono leggibili e con scale logaritmiche dove appropriato.
  \item Ogni grafico ha un ``purpose'' chiaro e viene commentato nel testo.
  \item Le fonti teoriche sono citate con precisione (numero di teorema/pagina).
  \item Le condizioni dei teoremi usati sono discusse (ipotesi soddisfatte o no?).
  \item Il confronto con un solutore di riferimento è presente per validazione.
  \item I parametri algoritmici sono stati studiati e le scelte sono motivate.
  \item Il report è stato inviato parzialmente ai docenti per feedback (consigliato!).
  \item Il questionario di valutazione del corso è stato compilato.
  \item Il repository \textbf{non} è pubblico.
\end{enumerate}

% ══════════════════════════════════════════════════════════════
\section{Timeline visuale}
\label{sec:timeline}

\begin{center}
\begin{tikzpicture}[x=1.4cm, y=0.8cm]
  % Grid
  \foreach \x in {0,...,10} {
    \draw[gray!30] (\x, -0.5) -- (\x, 6.5);
    \node[below, font=\footnotesize] at (\x, -0.5) {Sett.\ \x};
  }

  % Bars
  \fill[blue!30] (0, 5.5) rectangle (1, 6.2);
  \node[right, font=\small] at (0.05, 5.85) {Fase 0: Setup};

  \fill[green!30] (1, 4.5) rectangle (2, 5.2);
  \node[right, font=\small] at (1.05, 4.85) {Fase 1: ELM + Dati};

  \fill[orange!30] (2, 3.5) rectangle (4, 4.2);
  \node[right, font=\small] at (2.05, 3.85) {Fase 2: IRLS (A1)};

  \fill[red!30] (3, 2.5) rectangle (5, 3.2);
  \node[right, font=\small] at (3.05, 2.85) {Fase 3: Defl.\ Subgrad (A2)};

  \fill[purple!30] (5, 1.5) rectangle (7, 2.2);
  \node[right, font=\small] at (5.05, 1.85) {Fase 4: Esperimenti};

  \fill[cyan!30] (7, 0.5) rectangle (9, 1.2);
  \node[right, font=\small] at (7.05, 0.85) {Fase 5: Report};

  % Milestones
  \draw[thick, dashed, red] (4, -0.3) -- (4, 6.5);
  \node[above, font=\footnotesize\bfseries, red] at (4, 6.5) {Invio parziale};

  \draw[thick, dashed, red] (9, -0.3) -- (9, 6.5);
  \node[above, font=\footnotesize\bfseries, red] at (9, 6.5) {Consegna};
\end{tikzpicture}
\end{center}

\begin{tipbox}
Le fasi 2 e 3 sono parzialmente sovrapponibili: A e B possono lavorare in parallelo sui due algoritmi. Si consiglia di inviare il report parziale (sezione teorica + primi risultati di A1) dopo la settimana 4 per ricevere feedback.
\end{tipbox}

\end{document}
